% =====================================================================
% bridge-preamble.tex
% Preambule partage des livrables LaTeX du projet IoT Edge Bridge
% Auteur : Noel Junior Yando Fotso - EPHEC TS6 - 2025-2026
% Compilation : XeLaTeX + biber (latexmk -xelatex -shell-escape)
% =====================================================================

% --- Encodage et langue ----------------------------------------------
\usepackage[utf8]{inputenc}      % LuaTeX/XeTeX-friendly (ignore mais sans erreur)
\usepackage[T1]{fontenc}
\usepackage[french]{babel}
\usepackage{microtype}
\usepackage{csquotes}

% --- Fontes (XeLaTeX) ------------------------------------------------
% Detection conditionnelle : si XeLaTeX/LuaLaTeX, on essaie Fira Sans / Source Serif Pro,
% avec fallback automatique sur Latin Modern si la fonte est indisponible.
\usepackage{ifxetex,ifluatex}
\newif\ifxetexorluatex
\ifxetex\xetexorluatextrue\else
  \ifluatex\xetexorluatextrue\else\xetexorluatexfalse\fi
\fi

\ifxetexorluatex
  \usepackage{fontspec}
  % Source Serif Pro pour le rapport et le CDC ; Fira Sans pour le pitch.
  % Fallback silencieux sur Latin Modern si la fonte n'est pas installee.
  \IfFontExistsTF{Source Serif Pro}{%
    \setmainfont{Source Serif Pro}%
  }{%
    \IfFontExistsTF{Source Serif 4}{\setmainfont{Source Serif 4}}{}%
  }
  \IfFontExistsTF{Fira Sans}{%
    \setsansfont{Fira Sans}%
  }{}
  \IfFontExistsTF{Fira Mono}{%
    \setmonofont{Fira Mono}%
  }{}
\else
  % pdfLaTeX fallback : Latin Modern (par defaut, rien a faire).
\fi

% --- Graphiques et couleurs ------------------------------------------
\usepackage{graphicx}
\usepackage[table,dvipsnames,svgnames]{xcolor}

% Palette projet (charte)
\definecolor{bridgeblue}{HTML}{1F3A68}     % bleu profond, titres et cadres
\definecolor{bridgeteal}{HTML}{0FA4AF}     % teal d'accent
\definecolor{bridgegray}{HTML}{4A4A4A}     % gris texte secondaire
\definecolor{bridgealert}{HTML}{C62828}    % rouge alerte
\definecolor{bridgesoft}{HTML}{F2F6FB}     % bleu tres pale, fond doux
\definecolor{bridgeinfo}{HTML}{E3F4F4}     % teal tres pale, info
\definecolor{bridgewarn}{HTML}{FFF4E5}     % ambre tres pale, attention
\definecolor{bridgegreen}{HTML}{2E7D32}    % vert validation
\definecolor{bridgegreensoft}{HTML}{E8F5E9}% vert tres pale

% --- Liens hypertexte ------------------------------------------------
\usepackage[
  colorlinks=true,
  linkcolor=bridgeblue,
  citecolor=bridgeteal,
  urlcolor=bridgeteal,
  pdfborder={0 0 0}
]{hyperref}

% --- Tableaux et listes ----------------------------------------------
\usepackage{booktabs}
\usepackage{longtable}
\usepackage{tabularx}
\usepackage{array}
\usepackage{colortbl}
\usepackage{makecell}
\usepackage{enumitem}

% Style global de tableaux : zebrures bleu pale, regle haute teal, regles mid claires
\setlength{\arrayrulewidth}{0.4pt}
\renewcommand{\arraystretch}{1.25}
\arrayrulecolor{bridgeblue!40}

% Commande utilitaire pour activer les zebrures sur un tableau
% Usage : \begin{table} \bridgezebra ... \begin{tabular}...
\newcommand{\bridgezebra}{\rowcolors{2}{bridgesoft}{white}}

% Commande pour styliser une ligne d'en-tete (a placer dans la 1re ligne)
% Usage : \bridgeheader \textbf{Col1} & \textbf{Col2} \\ \midrule
\newcommand{\bridgeheader}{\rowcolor{bridgeblue!90}\color{white}\bfseries}

% Activation globale des zebrures par defaut sur tous les tableaux suivants.
% Pour desactiver localement : \rowcolors{1}{}{}
\rowcolors{2}{bridgesoft}{white}

% Style booktabs renforce avec couleur de regle
\let\oldtoprule\toprule
\let\oldmidrule\midrule
\let\oldbottomrule\bottomrule
\renewcommand{\toprule}{\arrayrulecolor{bridgeblue}\specialrule{0.7pt}{0pt}{0pt}\arrayrulecolor{bridgeblue!40}}
\renewcommand{\midrule}{\specialrule{0.4pt}{0pt}{0pt}}
\renewcommand{\bottomrule}{\arrayrulecolor{bridgeblue}\specialrule{0.7pt}{0pt}{0pt}\arrayrulecolor{bridgeblue!40}}

% Type de colonne avec en-tete bold blanc sur fond bleu
\newcolumntype{H}{>{\columncolor{bridgeblue}\color{white}\bfseries}l}

% Geometrie raisonnable (sera surchargee par les classes si besoin)
\usepackage{geometry}
\geometry{a4paper, margin=2.5cm, headheight=14pt}

% --- Boites informatives (tcolorbox) ---------------------------------
\usepackage[most]{tcolorbox}
\tcbset{
  bridgebox/.style={
    enhanced,
    colback=bridgesoft,
    colframe=bridgeblue,
    boxrule=0.4pt,
    arc=2pt,
    left=8pt,right=8pt,top=6pt,bottom=6pt,
    fonttitle=\bfseries\color{white},
    coltitle=white,
    colbacktitle=bridgeblue,
    title filled,
    breakable
  },
  bridgeinfobox/.style={
    bridgebox,
    colback=bridgeinfo,
    colframe=bridgeteal,
    colbacktitle=bridgeteal
  },
  bridgewarnbox/.style={
    bridgebox,
    colback=bridgewarn,
    colframe=orange!80!black,
    colbacktitle=orange!80!black
  },
  bridgekeybox/.style={
    bridgebox,
    colback=bridgegreensoft,
    colframe=bridgegreen,
    colbacktitle=bridgegreen
  }
}

% Macros raccourcies. title={#1} pour proteger les virgules dans le titre.
\newtcolorbox{bridgenote}[1][Note]{bridgebox, title={#1}}
\newtcolorbox{bridgeinfo}[1][Information]{bridgeinfobox, title={#1}}
\newtcolorbox{bridgewarn}[1][Attention]{bridgewarnbox, title={#1}}
\newtcolorbox{bridgekey}[1][A retenir]{bridgekeybox, title={#1}}

% --- Typographie KOMA-Script ----------------------------------------
% S'applique uniquement si on est en classe scrreprt/scrartcl/scrbook.
\makeatletter
\@ifclassloaded{scrreprt}{%
  \addtokomafont{disposition}{\color{bridgeblue}}%
  \addtokomafont{chapter}{\color{bridgeblue}\Huge\bfseries}%
  \addtokomafont{section}{\color{bridgeblue}}%
  \addtokomafont{subsection}{\color{bridgeblue}}%
  \addtokomafont{descriptionlabel}{\color{bridgeblue}\bfseries}%
  \KOMAoptions{numbers=noenddot, captions=tableheading, headings=optiontoheadandtoc}%
}{}
\@ifclassloaded{scrartcl}{%
  \addtokomafont{disposition}{\color{bridgeblue}}%
}{}
\makeatother

% --- Code et listings ------------------------------------------------
\usepackage{listings}

% Style generique
\lstdefinestyle{bridgebase}{%
  basicstyle=\ttfamily\small,
  keywordstyle=\color{bridgeblue}\bfseries,
  commentstyle=\color{bridgegray}\itshape,
  stringstyle=\color{bridgeteal},
  numberstyle=\tiny\color{bridgegray},
  numbers=left,
  numbersep=8pt,
  frame=single,
  rulecolor=\color{bridgegray!40},
  breaklines=true,
  showstringspaces=false,
  tabsize=2,
  captionpos=b
}

% Style JSON
\lstdefinelanguage{json}{%
  string=[s]{"}{"},
  stringstyle=\color{bridgeteal},
  comment=[l]{//},
  commentstyle=\color{bridgegray}\itshape,
  morestring=[b]",
  literate=
    *{0}{{{\color{bridgeblue}0}}}{1}
     {1}{{{\color{bridgeblue}1}}}{1}
     {2}{{{\color{bridgeblue}2}}}{1}
     {3}{{{\color{bridgeblue}3}}}{1}
     {4}{{{\color{bridgeblue}4}}}{1}
     {5}{{{\color{bridgeblue}5}}}{1}
     {6}{{{\color{bridgeblue}6}}}{1}
     {7}{{{\color{bridgeblue}7}}}{1}
     {8}{{{\color{bridgeblue}8}}}{1}
     {9}{{{\color{bridgeblue}9}}}{1}
     {:}{{{\color{bridgegray}{:}}}}{1}
     {,}{{{\color{bridgegray}{,}}}}{1}
     {\{}{{{\color{bridgeblue}{\{}}}}{1}
     {\}}{{{\color{bridgeblue}{\}}}}}{1}
     {[}{{{\color{bridgeblue}{[}}}}{1}
     {]}{{{\color{bridgeblue}{]}}}}{1},
}
\lstdefinestyle{bridgejson}{style=bridgebase, language=json}

% Style Python (utilise le langage builtin de listings)
\lstdefinestyle{bridgepython}{%
  style=bridgebase,
  language=Python,
  morekeywords={async,await,match,case}
}

\lstset{style=bridgebase}

% --- TikZ -------------------------------------------------------------
\usepackage{tikz}
\usetikzlibrary{shapes.geometric, arrows.meta, positioning, fit, backgrounds, calc}

% Styles partages pour les schemas du projet
\tikzset{
  bridgeblock/.style={%
    rectangle, rounded corners=2pt,
    draw=bridgeblue, thick,
    fill=bridgeblue!8,
    text=bridgeblue,
    align=center,
    minimum width=28mm, minimum height=10mm,
    inner sep=4pt,
    font=\small\sffamily
  },
  bridgeflow/.style={%
    -{Stealth[length=3mm,width=2mm]},
    draw=bridgeteal, thick
  },
  bridgealertstyle/.style={%
    rectangle, rounded corners=2pt,
    draw=bridgealert, thick,
    fill=bridgealert!8,
    text=bridgealert,
    align=center,
    minimum width=28mm, minimum height=10mm,
    font=\small\sffamily\bfseries
  },
  bridgenotetikz/.style={%
    rectangle, rounded corners=2pt,
    draw=bridgegray, dashed,
    fill=bridgegray!5,
    text=bridgegray,
    align=left,
    inner sep=4pt,
    font=\footnotesize\sffamily\itshape
  }
}

% Alias plus courts pour rester compatibles avec les schemas
\tikzset{bridgealert/.style={bridgealertstyle}}

% --- Bibliographie ---------------------------------------------------
\usepackage[
  backend=biber,
  style=ieee,
  sorting=nyt,
  maxbibnames=99,
  giveninits=true
]{biblatex}

% --- Sous-titre partage pour la page de garde ------------------------
% Les classes scrreprt et beamer ont leur propre logique de \subtitle ;
% on garde une commande explicite \bridgesubtitle pour compatibilite.
\providecommand{\bridgesubtitle}[1]{\gdef\@bridgesubtitle{#1}}
\providecommand{\thebridgesubtitle}{\@bridgesubtitle}

% Si la classe definit \subtitle (scrreprt KOMA, beamer), on l'utilise aussi.
\providecommand{\subtitle}[1]{\bridgesubtitle{#1}}

% --- Metadonnees PDF (renseignees a la compilation) ------------------
\AtBeginDocument{%
  \hypersetup{%
    pdfauthor={Noel Junior Yando Fotso},
    pdfsubject={IoT Edge Bridge - EPHEC TS6 2025-2026},
    pdfkeywords={IoT, FHIR, HL7, KMEHR, BIHR, MDR, RGPD, eSante, EPHEC}
  }
}
